\documentclass[10pt,letterpaper]{article}
\usepackage{technical-design-spec}
\setDocTitle{Real Time Character Animation Architecture}
\setDocSubtitle{Animation architecture and verification specification}
\setDocVersion{1.0}
\setDocStatus{Educational / Draft}
\setDocOwner{Game Development}
\setDocDate{2026-10-30}
\setDocSummary{Animation architecture and verification specification}
\begin{document}
\makedoctitle
\makedocsummary
\tableofcontents\clearpage

\section{Architecture stakeholders and concerns}
Gameplay design: responsive control and expressive intent. Animation: rig fidelity, contact, natural arcs. Engine: bounded time and allocation cost. QA: deterministic trace and stable regression results. Accessibility: readability at reduced motion and small displays.
\section{System context and logical interfaces}
\begin{verbatim}
Input -> Player Ability Rules -> Animation State Graph
                          -> Clip Sampler -> Pose Blender
Physics Contact Solver -> Contact Reconciler -> Final Pose
World Clock -> Fixed-Step Tick -> State Graph
Final Pose -> Renderer + Animation Event Dispatch
Telemetry <- each subsystem
\end{verbatim}
\section{Data ownership}
The authoritative simulation clock increments by integer ticks. Clip player owns normalized authored time; state machine owns selected state and interruption policy; physics owns world collision resolution; movement reconciler owns conversion between animation root delta and actor displacement. Renderer never advances gameplay state on its own.
\section{Functional requirements}
\Req{ANIM-R-001}{The runtime shall represent frame rate as a rational pair or an equivalent drift-free clock.}
\Req{ANIM-R-002}{The transition graph shall apply explicitly ordered guards and priorities, without implicit first-match dependence on container iteration.}
\Req{ANIM-R-003}{The root-motion integration shall have exactly one declared authoritative displacement owner.}
\Req{ANIM-R-004}{Animation markers shall have a documented boundary and loop-wrap delivery policy.}
\Req{ANIM-R-005}{A captured run shall include shot/clip IDs, fixed ticks, transition triggers and content versions.}
\Req{ANIM-R-006}{The runtime shall degrade safely under a fixed-step overload rather than skipping updates without telemetry.}
\section{Nonfunctional thresholds}
Populate platform-specific targets using the engineering-quality scorecard; record p50/p95 animation update time, CPU budget, memory/rig, contact drift, visual cue latency, continuity at clip boundaries, and accessibility configuration. This specification intentionally does not invent universal tolerances.
\section{Risk assessment}
High: double root displacement, accidental interrupt deadlocks, event duplication on looping. Medium: foot skating, retarget issues, asset mismatch. Observe through traces; mitigate with invariant assertions, integration tests, and owner sign-off.
\section{Traceability and validation}
\begin{tabular}{p{.21\linewidth}p{.33\linewidth}p{.3\linewidth}}
\textbf{Requirement}&\textbf{Evidence}&\textbf{Validation}\\\hline
ANIM-R-001&Frame rate field and frame boundary trace&24 vs $24000/1001$ conversion test\\
ANIM-R-002&Transition table&Interruptible attack state tests\\
ANIM-R-003&Root ownership record&Contact drift + actor displacement checks\\
ANIM-R-004&Marker semantics&Wrap, reverse, clamp test matrix\\
ANIM-R-005&Serialized trace fixture&Reproducibility replay check\\
ANIM-R-006&Overload monitor&Catch-up budget stress test\\
\end{tabular}
\section{Related source}
Chapters 8 and 73 identify the challenges of interactive-game animation; the software decomposition above is a supplementary engineering design.
\end{document}
